\section{关键实验指标}
\subsection{Iteration 与 win rate 观察}
Weston 把 self-rewarding 训练分作多个迭代，每一轮都采用 Self-Instruct 任务、verified response 和 judge 的组合。实验显示，第一轮训练就能让模型在 Alpaca Eval 2 上获得可观 win rate，而加入 meta judge 后 win rate 进一步提升，避免了原本在第三轮出现的 plateau（SRT 1908-2027、2784-2790）。

\begin{table}[H]
\centering
\begin{tabular}{p{4cm}p{4cm}p{7cm}}
\toprule
迭代阶段 & 样本量估算 & 观察 \\
\midrule
Seed Iteration & 少量 verified prompts & 模型快速掌握题型，但 reasoning 质量仍需 judge。 \\
Self-Rewarding Iteration 1 & Self-Instruct + verified response & Win rate 在 Alpaca Eval 2 初步提升，模型学会自我评分。 \\
Self-Rewarding Iteration 2 & 新任务 + cross-check & 反复验证错误，reward 向正确方向偏移。 \\
Meta-Judge Iteration & 加入 SelfT evals judge & 解决 plateau，verified response 的可信度变得可量化，win rate 稳步上升。 \\
\bottomrule
\end{tabular}
\caption{讲座中展示的 iteration 与 win rate 观察（qualitative）。}
\end{table}

\subsection{SelfT evaluators 及其他数据源}
除了 Alpaca Eval 2，Weston 还提到 SelfT、Self-Feeding Chatbot 以及 generic instruction data（数学/代码/判断/对话）。这些数据源相互补充：SelfT 提供 judge supervision，Self-Feeding 提供 conversational feedback，而 Self-Instruct 提供 new tasks，形成 ``模型自训练'' 所需的四角支撑。

\begin{importantbox}{数据源四角}
\begin{itemize}
    \item Self-Instruct：生成 instruction → input → output。
    \item Self-Feeding Chatbot：通过对话生成 verified response。
    \item SelfT evaluators：提供 judge 仿真标签。
    \item Alpaca Eval 2：benchmark 上的 win rate 观测，用于指导 reward 模型。
\end{itemize}
\end{importantbox}

\subsection{本章小结}
不同迭代在数据源和 judge 上逐步加码：Self-Instruct 负责任务、Self-Feeding 提供验证、SelfT 训练 judge、Alpaca Eval 2 提供最终 win rate，四者配合稳步提升 reasoning 表现。

